\chapter{Compactified Inverse-Boundary Geometry and the Unique Half-Layer}
\label{ch:compactified-inverse-boundary}

\status{Exact global geometric and quotient reduction; quotient real-rootedness remains open.}

\section{Purpose}

This chapter replaces an unnecessarily local proof emphasis with the global projective geometry already present in the inverse-boundary coordinate. No Hermite finite-section assumption is used here.

Define
\[
\Omega(s)=\frac{s}{1-s}.
\]
On the Riemann sphere,
\[
\Omega(0)=0,
\qquad
\Omega(1)=\infty.
\]
Thus the two real strip boundaries become the reciprocal projective pair $0\leftrightarrow\infty$.

For the critical-axis reflection
\[
K(s)=1-\overline{s},
\]
we have
\[
\Omega(K(s))=\frac{1}{\overline{\Omega(s)}}.
\]
Therefore, for
\[
R(s):=|\Omega(s)|,
\]
the induced radial involution is exactly
\[
R\longmapsto R^{-1}.
\]

\section{Compactifying zero and infinity}

Define the compactified radial coordinate
\[
\boxed{
q(s):=\frac{R(s)}{1+R(s)}
}
\]
with the extended values $q=0$ at $R=0$ and $q=1$ at $R=\infty$. Hence
\[
[0,\infty]\longrightarrow[0,1],
\qquad
0\mapsto0,
\qquad
\infty\mapsto1.
\]

Under reciprocal inversion,
\[
q\left(R^{-1}\right)
=\frac{R^{-1}}{1+R^{-1}}
=\frac{1}{1+R}
=1-q(R).
\]
Thus the projective inversion becomes the ordinary complement involution
\[
\boxed{q\mapsto1-q}.
\]

\begin{theorem}[Compactified inverse-boundary half-axis theorem]
The Möbius coordinate $\Omega(s)=s/(1-s)$ followed by radial compactification $q=|\Omega|/(1+|\Omega|)$ sends the reciprocal boundary pair $0\leftrightarrow\infty$ to $0\leftrightarrow1$, conjugates radial inversion to $q\mapsto1-q$, and sends the critical line $\Re s=1/2$ exactly to the unique self-dual layer $q=1/2$.
\end{theorem}

\begin{proof}
The complement involution $q\mapsto1-q$ has the unique fixed point $q=1/2$. Moreover,
\[
q=\frac12
\iff
\frac{R}{1+R}=\frac12
\iff
R=1.
\]
By \cref{chap:li-negative-inverse},
\[
R=|\Omega(s)|=1
\iff
\Re s=\frac12.
\]
Therefore the unique compactified self-dual radial layer is precisely the critical line.
\end{proof}

\section{Centred signed radius}

Define
\[
\boxed{
\eta(s):=2q(s)-1=\frac{R(s)-1}{R(s)+1}
}.
\]
Then
\[
R=0\mapsto\eta=-1,
\qquad
R=1\mapsto\eta=0,
\qquad
R=\infty\mapsto\eta=+1.
\]
Reciprocal inversion gives
\[
\boxed{\eta\mapsto-\eta}.
\]

If
\[
B(s)=\log R(s)=\log|\Omega(s)|,
\]
then
\[
\eta(s)
=\frac{e^{B(s)}-1}{e^{B(s)}+1}
=\boxed{\tanh\!\left(\frac{B(s)}{2}\right)}.
\]
Consequently,
\[
\boxed{
B=0
\iff
\eta=0
\iff
q=\frac12
\iff
\Re s=\frac12.
}
\]

\section{Completed zeta and the radial bridge}

Set
\[
X(u):=\xi\!\left(\frac{u}{1+u}\right),
\qquad u=\Omega(s).
\]
The Riemann Hypothesis is exactly equivalent to
\[
X(u)=0\Longrightarrow |u|=1,
\]
and therefore to
\[
\boxed{X(u)=0\Longrightarrow q(u)=\frac12.}
\]

\begin{problem}[SOH-G001: global self-dual zero closure]
Prove
\[
X(u)=0\Longrightarrow q(u)=\frac12.
\]
\end{problem}

\section{Holomorphic quotient of the inversion}

The radial coordinate $q$ is geometrically canonical but not holomorphic. For the analytic zero problem put
\[
z=s-\frac12,
\qquad
\Xi_c(z)=\xi\!\left(\frac12+z\right).
\]
The functional equation gives $\Xi_c(z)=\Xi_c(-z)$. Hence $\Xi_c$ is even entire and there is a unique entire function $F$ such that
\[
\boxed{\Xi_c(z)=F(z^2).}
\]
Writing $u=\Omega(s)$ gives
\[
z=\frac{u-1}{2(u+1)},
\qquad
\boxed{w:=z^2=\frac{(u-1)^2}{4(u+1)^2}.}
\]
A direct calculation shows
\[
w(1/u)=w(u),
\]
so $w$ exactly quotients the functional inversion $u\leftrightarrow1/u$.

For $|u|=1$, write $u=e^{i\theta}$. Then
\[
\frac{u-1}{u+1}=i\tan\frac{\theta}{2},
\]
therefore
\[
\boxed{w=-\frac14\tan^2\frac{\theta}{2}\in(-\infty,0].}
\]
Conversely $w=z^2\le0$ implies $z\in i\mathbb R$, hence $\Re s=1/2$. Thus
\[
\boxed{
q=\frac12
\iff |u|=1
\iff w\in(-\infty,0].
}
\]

\section{Coefficient positivity removes the positive real axis}

Write
\[
F(w)=\sum_{n=0}^{\infty}a_n w^n.
\]
The positivity of the non-zero Taylor coefficients of $\xi$ about $s=1/2$ is established independently by integral formulas in \citet{defranco2019xi}. In this quotient normalization this gives $a_n>0$ for the non-zero coefficients. Hence
\[
\boxed{F(x)>0\qquad(x\ge0).}
\]
Therefore $F$ has no zero on the non-negative real axis independently of RH.

It follows that the remaining target for this quotient route is sharper than SOH-G001:

\begin{problem}[SOH-G003: quotient real-zero target]
Prove that every zero of the entire quotient function $F$ is real.
\end{problem}

If SOH-G003 holds, coefficient positivity forces every zero of $F$ onto the negative real half-axis, which by the exact quotient geometry is equivalent to the critical line. Conversely RH implies that every quotient zero lies on that half-axis. Hence
\[
\boxed{
\mathrm{RH}
\iff
\text{all zeros of }F\text{ are real}.
}
\]

\gap{The projective geometry, the inversion quotient, and exclusion of the positive real half-axis are closed. Within this quotient route, the remaining RH-equivalent analytic statement is quotient real-rootedness. Other branches of the monograph retain their own separately stated open conditions.}

\section{Proof-target firewall}

The following statements are exact and must not be reopened as separate numerical subproblems:
\begin{enumerate}
\item the projective boundary pair is $0\leftrightarrow\infty$;
\item radial inversion is $R\mapsto R^{-1}$;
\item compactification converts inversion to $q\mapsto1-q$;
\item the unique self-dual layer is $q=1/2$;
\item that layer is exactly $\Re s=1/2$;
\item $w=(u-1)^2/[4(u+1)^2]$ holomorphically quotients $u\leftrightarrow1/u$;
\item the self-dual circle maps exactly to $w\le0$;
\item positivity of the centered Taylor coefficients excludes $w\ge0$ from the zero set.
\end{enumerate}

Finite-Hermite, moving-boundary, and prime-tail constructions may remain as auxiliary diagnostics, but they are not independent prerequisites for SOH-G003 and must not displace the global quotient target.

\begin{warning}
This chapter is an exact reduction, not a proof of the Riemann Hypothesis. SOH-G003 remains open until a non-circular analytic argument proves real-rootedness of the quotient entire function $F$.
\end{warning}
